\subsection{CANONICAL DECOMPOSITION OF A CONTINUOUS MAPPING}

Let X and Y be two topological spaces, \(f: X \rightarrow Y\) a continuous map, R the equivalence relation \(f(x) = f(y)\) on X. Consider the canonical decomposition

\[
f: \mathrm{X} \xrightarrow {\varphi} \mathrm{X} / \mathrm{R} \xrightarrow {g} f (\mathrm{X}) \xrightarrow {\psi} \mathrm{Y}
\]

where \(\varphi\) is the canonical (surjective) mapping of X onto the quotient space X/R, \(\psi\) is the canonical injection of the subspace \(f(X)\) into Y, and \(g\) is the bijection associated with \(f\) (Set Theory R, § 5, no. 3). It is immediately seen that \(g\) is continuous (by Proposition 6 of no. 4); this is also a particular case of a general result on quotient structures. (Cf. Set Theory, Chapter IV, § 2, no. 6). But the bijection \(g\) is not necessarily a homeomorphism.

\pandocbounded{\includegraphics[keepaspectratio]{images/part001_a90573a3.jpg}}

PROPOSITION 8. Let \(f = \psi \circ g \circ \varphi\) be the canonical decomposition of a continuous mapping \(f: X \to Y\) , and let \(R\) denote the equivalence relation

\[
f (x) = f (y).
\]

Then the following three conditions are equivalent:

\begin{enumerate}
\def\labelenumi{\alph{enumi})}
\item
  \(g\) is a homeomorphism of \(\mathbf{X} / \mathbf{R}\) onto \(f(\mathbf{X})\) .
\item
  The image under \(f\) of every open set which is saturated with respect to \(\mathbf{R}\) is an open set in the subspace \(f(\mathbf{X})\) .
\item
  The image under \(f\) of every closed set which is saturated with respect to \(\mathbf{R}\) is a closed set in the subspace \(f(\mathbf{X})\) .
\end{enumerate}

For the condition b) {[}resp. c){]} expresses that the image under g of every open (resp. closed) set in X/R is an open (resp. closed) set in \(f(\mathbf{X})\) .

Example. Let X be a topological space, \((\mathbf{X}_{i})_{i\in I}\) a covering of X, Y the sum of the subspaces \(X_{v}\) of X; then there is a partition \((\mathbf{Y}_{i})_{i\in I}\) of Y into subspaces which are both open and closed, and for each \(i\in I\) there is a homeomorphism \(f_{i}:Y_{i}\to X_{i}\) . Let \(f:Y\to X\) be the continuous mapping which agrees with \(f_{i}\) on \(Y_{i}\) for each \(i\in I\) , and let R be the equivalence relation \(f(x)=f(y)\) ; the quotient space Y/R is thus obtained by ``pasting together'' the \(Y_{i}\) (§2, no. 5). Consider the bijection \(g:Y/R\to X\) associated with f; in general g is not a homeomorphism, as is shown by the example in which each \(X_{i}\) consists of a single point and X is not discrete. However, if the interiors of the \(X_{i}\) cover X, or if \((\mathbf{X}_{i})\) is a locally finite closed covering of X, then g is a homeomorphism: for if U is any open subset of Y which is saturated with respect to R, then for each \(i\in I\) the set

\[
f (\mathbf {U}) \cap \mathbf {X} _ {i} = f _ {i} (\mathbf {U} \cap \mathbf {Y} _ {i})
\]

is open in \(X_{i}\) , and the assertion follows from Proposition 3 of no. 1.

The following proposition gives a simple sufficient condition for \(g\) to be a homeomorphism:

PROPOSITION 9. Let \(f: \mathbf{X} \to \mathbf{Y}\) be a continuous surjection, and let \(\mathbf{R}\) denote the equivalence relation \(f(x) = f(y)\) . If there is a continuous section \(s: \mathbf{Y} \to \mathbf{X}\) associated with \(f\) (Set Theory, Chapter II, § 3, no. 8, Definition 11), then the mapping \(g: \mathbf{X} / \mathbf{R} \to \mathbf{Y}\) associated with \(f\) is a homeomorphism, and \(s\) is a homeomorphism of \(\mathbf{Y}\) onto the subspace \(s(\mathbf{Y})\) of \(\mathbf{X}\) .

For if \(\varphi: \mathbf{X} \to \mathbf{X} / \mathbf{R}\) is the canonical mapping, then \(g\) and \(\varphi \circ s\) are bijective, continuous and inverse to each other; likewise \(s\) and the restriction of \(f\) to \(s(\mathbf{Y})\) are bijective, continuous and inverse to each other.

If R is an equivalence relation on a topological space X and

\[
\varphi : \mathbf {X} \rightarrow \mathbf {X} / \mathbf {R}
\]

is the canonical mapping, a continuous section \(s: \mathbf{X} / \mathbf{R} \to \mathbf{X}\) associated with \(\varphi\) is also called a continuous section of \(\mathbf{X}\) with respect to \(\mathbf{R}\) (cf.~Set Theory, Chapter II, § 6, no. 2); the subspace \(s(\mathbf{X} / \mathbf{R})\) of \(\mathbf{X}\) is then homeomorphic to \(\mathbf{X} / \mathbf{R}\) . If we are given \(s(\mathbf{X} / \mathbf{R})\) , then \(s\) is uniquely determined; \(s(\mathbf{X} / \mathbf{R})\) is often called, by abuse of language, a (continuous) section of \(\mathbf{X}\) with respect to \(\mathbf{R}\) .

\pandocbounded{\includegraphics[keepaspectratio]{images/part001_feb98675.jpg}}

A continuous section of a topological space with respect to an equivalence relation need not exist (Exercise 12).

